\documentclass[10pt,a4paper]{article}
\usepackage[utf8]{inputenc}
\usepackage[french]{babel}
\usepackage[T1]{fontenc}
\usepackage{amsmath}
\usepackage{amsfonts}
\usepackage{amssymb}
\usepackage{graphicx}
\usepackage{tikz}
\usepackage{jigsaw}
\usepackage{geometry}
\geometry{landscape,top=0.5cm,bottom=0.5cm,left=0.5cm,right=0.5cm}
\pagestyle{empty}
\usepackage{csvsimple}
\begin{document}
{\centering Associer les pièces par quatre comme dans l'exemple de gauche pour former six raisonnements (on n'écrira que les numéros des pièces).

Le texte d'une pièce peut être lu à l'envers.

Pour chaque raisonnement, tracer une figure à main levée.

}

\vspace{1cm}
\begin{tikzpicture}
\def\sc{3}
\def\xshift{-5}
\begin{scope}[xshift=\xshift cm,yshift=-\sc*0 cm,scale=\sc]
 \piece{-1}{0}{0}{0}
 \draw (0.5,0.5) node[text width=\sc cm,align=center] {

};
\end{scope}
\begin{scope}[xshift=\xshift cm,yshift=-\sc*1 cm,scale=\sc]
 \piece{-1}{0}{1}{0}
 \draw (0.5,0.5) node[text width=\sc cm,align=center] {

};
\end{scope}
\begin{scope}[xshift=\xshift cm,yshift=-\sc*2 cm,scale=\sc]
 \piece{1}{0}{1}{0}
 \draw (0.5,0.5) node[text width=\sc cm,align=center] {

};
\end{scope}
\begin{scope}[xshift=\xshift cm,yshift=-\sc*3 cm,scale=\sc]
 \piece{0}{0}{-1}{0}
 \draw (0.5,0.5) node[text width=\sc cm,align=center] {

};
\end{scope}
\def\sc{3}
\def\ecart{4}
\csvreader{./pieces.csv}{1=\numli,2=\numcol,3=\enbas,4=\adroite,5=\enhaut,6=\agauche,7=\phrase}{
\begin{scope}[xshift=\numcol*\ecart cm,yshift=-\numli*\ecart cm,scale=\sc]
 \piece{\enbas}{\adroite}{\enhaut}{\agauche}
 \draw (0.5,0.5) node[text width=\sc cm,align=center] {
\phrase
};
\draw (0.15,0.12) node {\Huge \thecsvrow};
\end{scope}
}
\end{tikzpicture}
\end{document}